\section{Intrinsic information in the enveloping algebra}\label{sec:intrinsic}
Throughout this section, $k$ has characteristic zero and all Lie algebras
are finite dimensional. We use PBW to regard a Lie algebra as a subspace
of its enveloping algebra. If $J$ is a Lie ideal of $g$, then
$U(g)J$ denotes the two-sided ideal it generates, and
\[
 U(g)/U(g)J\simeq U(g/J).
\]
The standard Lie-theoretic ingredients are PBW, Lie's theorem, Engel's
theorem, and the existence of the nilradical; see
\cite{Humphreys1972,BourbakiLie1989}.

\subsection{Solvability, augmentation, and scalar extension}

\begin{lemma}\label{lem:target-solvable}
If $L$ is solvable and
$U(L)\simeq U(M)$ as unital $k$-algebras, then
$M$ is solvable.
\end{lemma}
\begin{proof}
Over an algebraically closed field, a finite-dimensional Lie algebra is
solvable if and only if all its finite-dimensional irreducible modules
are one dimensional. The forward implication is Lie's theorem. For the
converse, a composition series of the adjoint module becomes a full
invariant flag. The adjoint image is therefore upper triangular and
solvable; its central kernel implies solvability of the original algebra.

For an algebraic closure $\Omega/k$, PBW gives
\[
 \Omega\otimes_k U(g)\simeq U(\Omega\otimes_kg).
\]
An associative algebra isomorphism transports simple modules without
changing their underlying vector spaces. Thus the preceding criterion
passes from $L_\Omega$ to $M_\Omega$. Derived series
commute with field extension, and faithful scalar extension reflects
vanishing, so solvability descends to $k$.
\end{proof}

Denote the initially given associative isomorphism by $\Phi_{\mathrm{raw}}$.
\begin{lemma}\label{lem:augmentation}
An isomorphism $\Phi_{\mathrm{raw}}:U(L)\to U(M)$ can be made
augmentation preserving by composing it with an explicit automorphism
of its source.
\end{lemma}
\begin{proof}
Let $\varepsilon_{g}$ denote the standard augmentation and set
$\chi=\varepsilon_{M}\Phi_{\mathrm{raw}}|_{L}$. Then
$\chi([x,y])=0$. For any Lie character $\lambda$, the assignment
$x\mapsto x+\lambda(x)1$ preserves the enveloping relations and extends
to an automorphism $\tau_\lambda$, with inverse $\tau_{-\lambda}$.
Consequently
\[
 \Phi=\Phi_{\mathrm{raw}}\circ\tau_{-\chi},\qquad
 \varepsilon_{M}\Phi=\varepsilon_{L}.
\]
It suffices to check the latter identity on the Lie generators.
\end{proof}
From now on, $\Phi$ denotes this specified normalized isomorphism;
$\Phi_{\mathrm{raw}}$ denotes the initially given one.

Write $N(g)$ for the nilradical and
$\gamma_1N=N$, $\gamma_{r+1}N=[N,\gamma_rN]$ for its lower central
series. We will need compatibility with arbitrary field extensions,
including transcendental ones. A finite system of trace equations
provides this directly. The use of the full associative algebra generated
by the adjoint operators is important: testing only the individual
adjoint operators would replace these equations by the Killing form,
which does not give the assertion below.

\begin{lemma}\label{lem:trace-nilradical}
Let $\mathcal A\subseteq\End_k(g)$ be the unital associative
algebra generated by $\ad(g)$. Then
\[
 N(g)=
 \{x\in g:\Tr(\ad(x)T)=0\text{ for every }T\in\mathcal A\}.
\]
\end{lemma}
\begin{proof}
Call the displayed right-hand side $J$. Cyclicity of trace gives
\[
 \Tr(\ad([y,x])T)
 =\Tr\bigl(\ad(x)(T\ad(y)-\ad(y)T)\bigr),
\]
so $J$ is an ideal. If $x\in J$, choosing
$T=(\ad x)^{r-1}$ shows that all positive power traces of $\ad x$
vanish. In dimension $n$, the first $n$ Newton identities, whose
integer denominators are invertible in $k$, make all the nonleading
coefficients of its characteristic polynomial zero. Cayley--Hamilton
then implies that $\ad x$ is nilpotent. Its restriction to $J$ is nilpotent, so Engel's theorem
makes $J$ nilpotent. Hence $J\subseteq N(g)$.

Conversely, let $I$ be a nilpotent ideal and $x\in I$. Every element
of $\mathcal A$ preserves every Lie ideal, in particular the ideals
$\gamma_rI$. Thus $\ad(x)T$ sends $g$ into $I$ and sends
$\gamma_rI$ into $\gamma_{r+1}I$. It is nilpotent and has trace zero.
Therefore $I\subseteq J$, proving the reverse inclusion.
\end{proof}

\begin{corollary}\label{cor:nilradical-extension}
For every field extension $K/k$,
\[
 N(g_K)=N(g)_K,\qquad
 \gamma_rN(g_K)=(\gamma_rN(g))_K.
\]
\end{corollary}
\begin{proof}
Choose a basis $T_1,\ldots,T_s$ of $\mathcal A$. The nilradical is
the kernel of the linear map
$x\mapsto(\Tr(\ad(x)T_i))_i$. Each extended $T_i$ belongs to the
adjoint envelope over $K$. By \cref{lem:trace-nilradical}, the new
nilradical satisfies these extended equations. Kernels commute with
field extension, giving $N(g_K)\subseteq N(g)_K$.
The reverse inclusion follows because extension preserves nilpotent
ideals. The lower central identities follow by extending the successive
images of the bracket maps.
\end{proof}

\subsection{Weights and two derivation lemmas}

\begin{lemma}\label{lem:weight-separation}
For a solvable $g$, put $N=N(g)$,
$Q=g/N$, and $V=N/[N,N]$. Then $Q$ is abelian. Over an
algebraic closure, the simultaneous weights of its action on $V$ span
$Q^*$; each such weight has a nonzero common eigenvector.
\end{lemma}
\begin{proof}
Lie's theorem makes $\ad(g)$ upper triangular over an
algebraic closure. Its commutators are strictly upper triangular, so
Engel's theorem shows that $[g,g]$ is nilpotent.
This descends to $k$, giving $[g,g]\subseteq N$.
The action on $V$ factors through the abelian quotient $Q$.

We first observe that a derivation $D$ of a nilpotent algebra $N$ is
nilpotent if its action on $V$ is nilpotent. Iterated brackets give
surjective $D$-equivariant maps
\[
 V^{\otimes r}\longrightarrow\gamma_rN/\gamma_{r+1}N.
\]
Replacing one argument by an element of $[N,N]$ changes the bracket
by an element of $\gamma_{r+1}N$, which justifies these maps. If
$D^s|_V=0$, its action on $V^{\otimes r}$ is a sum of $r$ commuting
operators and has power $r(s-1)+1$ equal to zero. Thus $D$ is nilpotent
on every successive quotient and hence on $N$.

Work now over the algebraic closure, using
\cref{cor:nilradical-extension}. Let $Q_0$ be the common kernel of all
weights and $H$ its inverse image in $g$. This is an ideal.
Every $x\in H$ acts nilpotently on $V$, therefore on $N$, and acts
trivially on $H/N$. Hence every $\ad_Hx$ is nilpotent. Engel's theorem
and maximality of $N$ imply $H=N$, so $Q_0=0$, proving separation.
On each simultaneous generalized weight space, the commuting
nilpotent parts have a nonzero common kernel. This supplies a genuine
common eigenvector for each weight.
\end{proof}

The proof also covers $Q=0$, when separation in the zero dual space
is automatic. In general, a finite subfamily of the weights already
separates $Q$, since $Q$ is finite dimensional.
Zero weights and nonsemisimple actions are allowed. Separation of the
weights, rather than mere faithfulness of the action, is essential below.

\begin{lemma}\label{lem:leading-LF}
Suppose $A$ has an exhaustive increasing algebra filtration $F_dA$,
and a locally finite derivation $D$ satisfies
$D(F_dA)\subseteq F_{d+r}A$ for a fixed $r>0$. Its induced homogeneous
derivation of degree $r$ on $\gr_FA$ is locally nilpotent.
\end{lemma}
\begin{proof}
Lift a homogeneous symbol of degree $d$ to $b\in F_dA$. The
finite-dimensional orbit span of $b$ lies in some $F_BA$. The $j$th
iterate of the induced operator is the symbol of $D^jb$ in degree
$d+jr$, hence is zero once $d+jr>B$. Finite sums of homogeneous
symbols are handled by taking a maximum. The individual $F_dA$ need
not be finite dimensional.
\end{proof}

\begin{lemma}\label{lem:LND-divisibility}
If $D$ is a locally nilpotent derivation of a commutative domain of
characteristic zero and $0\ne u$ satisfies $D(u)\in uA$, then $D(u)=0$.
\end{lemma}
\begin{proof}
For $a\ne0$ set $\nu(a)=\max\{j:D^ja\ne0\}$. The iterated Leibniz
formula gives $\nu(ab)=\nu(a)+\nu(b)$: the unique surviving term at
that order has a nonzero binomial coefficient:
\[
 D^{\nu(a)+\nu(b)}(ab)
 =\binom{\nu(a)+\nu(b)}{\nu(a)}D^{\nu(a)}a\,D^{\nu(b)}b.
\]
If $D(u)=uv\ne0$, then
$\nu(u)-1=\nu(Du)=\nu(u)+\nu(v)$, a contradiction.
\end{proof}

\subsection{The abelian-extension calculation}
For an associative algebra $A$, let $\LN(A)$ and $\LF(A)$ be the
sets of elements whose inner derivations are respectively locally
nilpotent and locally finite. Local finiteness means that each vector
has a finite-dimensional span of iterates. These definitions are
intrinsic to the associative algebra.

\begin{proposition}\label{prop:LF-LN}
Let $H$ be an abelian ideal of $g$ with abelian quotient
$T=g/H$. Suppose the simultaneous weights of $T$ on $H$
after algebraic closure span $T^*$. Then, inside $U(g)$,
\[
 \LN(U(g))=S(H),\qquad
 \LF(U(g))=g+S(H)=s(T)+S(H)
\]
for every linear section $s:T\to g$.
\end{proposition}
\begin{proof}
First work over an algebraically closed field. Choose lifts
$t_1,\ldots,t_d$ of a basis of $T$, and put
$\delta_i=\ad(t_i)|_{S(H)}$. These are commuting linear polynomial
derivations: $[\delta_i,\delta_j]=\ad([t_i,t_j])|_{S(H)}=0$.
We retain the possibly nonzero cocycle $c_{ij}=[t_i,t_j]\in H$.

Filter by the number of quotient letters: $H$ has degree zero and
each $t_i$ degree one, and we set $F_{-1}=0$. PBW gives
\[
 \gr U(g)=S(H)[\xi_1,\ldots,\xi_d].
\]
Indeed $[t_i,H]\subseteq H$ lowers degree by one and $[t_i,t_j]\in H$
lowers degree by two. In particular $[F_m,F_n]\subseteq F_{m+n-1}$.

Let $a\in\LF(U(g))$ have degree $m\ge2$ and leading
symbol $p$. The derivation $\ad a$ induces a derivation $D$ of
positive degree $m-1$, locally nilpotent by \cref{lem:leading-LF}.
For a common eigenvector $0\ne w\in H$ of weight $\lambda$, the
commutator Leibniz rule yields
\[
 D(w)=w\sum_i\lambda(t_i)\frac{\partial p}{\partial\xi_i}.
\]
Here each occurrence of $t_i$ contributes $\lambda(t_i)$; further
reorderings lower degree. By \cref{lem:LND-divisibility}, the sum is
zero. Separation by the weights forces every partial derivative of
$p$ to vanish. Characteristic zero makes $p$ independent of the
$\xi_i$, contradicting its positive homogeneous degree. Thus $m\le1$.

The cocycle has not been set to zero: its contributions to $[a,t_j]$
have degree at most $m-1$, whereas the symbol defining $D(t_j)$ is
taken in degree $m$. Terms from commuting coefficients in $S(H)$
past $t_j$ may survive in degree $m$ and are retained. This verifies
the degree estimate also on the quotient generators.

Write, by PBW,
\[
 a=\sum_i a_i(H)t_i+b(H).
\]
On the invariant subalgebra $S(H)$, its inner derivation is the
locally finite derivation $\sum_i a_i\delta_i$. If the largest
ordinary polynomial degree of the coefficients $a_i$ is $e>0$, its
leading derivation is $E=\sum_i(a_i)_e\delta_i$, of positive degree
$e$. Again it is locally nilpotent. For each eigenvector as above,
\[
 E(w)=w\sum_i\lambda(t_i)(a_i)_e.
\]
Divisibility and separation force all $(a_i)_e$ to vanish, a
contradiction. Thus every $a_i$ is constant, proving the required
inclusion for $\LF$.

Conversely take $a=s(t)+b(H)$ and give the basis elements of $H$
weight one and the lifts weight $e=\max(1,\deg b)$, using $e=1$
when $b=0$. Bounded weighted PBW spaces are finite dimensional.
The relations have nonincreasing weight, and
\[
 [a,H]\subseteq H,\qquad
 [a,s(q)]=[s(t),s(q)]-\delta_qb
\]
has polynomial degree at most $e$. Thus $\ad a$ preserves every
such space and is locally finite. This proves the $\LF$ equality
with the cocycle included.

For $b\in S(H)$, $\ad b$ vanishes on $S(H)$ and strictly lowers
the number of quotient letters; it is locally nilpotent. Conversely
$\LN\subseteq\LF$, so an element of $\LN$ is $s(t)+b(H)$. If
$t\ne0$, some weight has $\lambda(t)\ne0$, and on its eigenvector
\[
 (\ad(s(t)+b))^rw=\lambda(t)^rw\ne0
\]
for every $r$. Hence $t=0$.

Finally, the two local properties of an operator defined over $k$
can be checked after field extension. In the forward direction each
extended vector is a finite sum of pure tensors: take the maximum
of finitely many nilpotence exponents, or the sum of finitely many
finite-dimensional orbit spaces. Conversely, an infinite independent
family of iterates remains independent after extension, and a nonzero
iterate remains nonzero. These assertions concern an operator defined
over $k$ and do not require any closure property of the sets $\LN$
and $\LF$ in an arbitrary associative algebra.
The memberships in $S(H)$ and $s(T)+S(H)$ are vanishing conditions
on PBW coefficients and therefore descend. This proves the assertion
over the original field. The equality
$g+S(H)=s(T)+S(H)$ follows from
$g=H+s(T)$.
\end{proof}

\begin{remark}
The extension in \cref{prop:LF-LN} need not split. For example, let
$g$ have basis $t_1,t_2,w_1,w_2,z$ and only the nonzero
brackets
\[
 [t_1,w_1]=w_1,\qquad [t_2,w_2]=w_2,\qquad [t_1,t_2]=z.
\]
Then $H=\langle w_1,w_2,z\rangle$ is abelian and its two nonzero
weights separate $g/H$. Changing the lifts adds only
$w_1,w_2$ components to their bracket, so its $z$ component cannot
be removed. The proposition applies to this extension without change.
\end{remark}

\begin{proposition}\label{prop:intrinsic-ideal}
For a solvable $g$ with nilradical $N$,
\[
 U(g)[N,N]
 =\bigl\langle[a,b]:a,b\in\LN(U(g))\bigr\rangle_{\mathrm{ideal}}.
\]
Consequently every associative enveloping-algebra isomorphism preserves
this ideal, without an augmentation hypothesis.
\end{proposition}
\begin{proof}
For $x\in N$, $(\ad x)^r(g)\subseteq\gamma_rN$ for
$r\ge1$. Hence $\ad x$ is nilpotent on the Lie generators and
locally nilpotent on their products, by the iterated Leibniz rule.
Thus $N\subseteq\LN(U(g))$, giving one inclusion.

In the reverse direction, set $\bar{g}=g/[N,N]$.
Its abelian ideal $V=N/[N,N]$ and quotient $Q=g/N$
satisfy \cref{lem:weight-separation}. By \cref{prop:LF-LN},
$\LN(U(\bar{g}))=S(V)$ is commutative. Locally nilpotent
inner derivations descend through associative quotients, so all the
displayed commutators vanish in $U(\bar{g})$.
Finally, an algebra isomorphism preserves $\LN$ and commutators.
\end{proof}

We have therefore recovered an actual quotient enveloping algebra and,
inside it, an actual polynomial subalgebra through intrinsic associative
conditions. For two isomorphic enveloping algebras of solvable Lie
algebras, the induced quotient isomorphism carries $S(N/[N,N])$ onto
the corresponding polynomial algebra on the other side, and preserves
the set of linear Lie elements plus that polynomial algebra. These are
the inputs for recovering the filtration in the next section; preservation
of the original subalgebra $U(N)$ has not been asserted.
